\begin{lemma}
Let $\Omega\subseteq\mathbb{C}$ be open, let $Z\subseteq\Omega$ meet every
compact subset of $\Omega$ in a finite set, let $k\ge1$ be an integer and
$\alpha>0$, and let $w\colon\mathbb{C}\to\mathbb{C}$ satisfy
\begin{enumerate}
\item[(H1)] $w$ is of class $C^1$ on $\Omega$;
\item[(H2)] $w$ is of class $C^\infty$ on $\Omega-Z$;
\item[(H3)] for every $z_0\in Z$ there are $\delta>0$ and constants $B_j$ with
\[
  \sum_{|v|=j}\bigl|D^{v}w(z)\bigr|\le B_j\,|z-z_0|^{\,k+\alpha-j}
  \qquad(1\le j\le k,\ 0<|z-z_0|<\delta),
\]
the sum running over the words $v$ of length $j$ in $\{0,1\}$ and $D^{v}$ being
the iterated partial derivative of \noderef{PartialWord}.
\end{enumerate}
Then $w$ is of class $C^{k}$ on $\Omega$.

This is the step Section 6 of the paper passes over.  Under the hypotheses of
Theorem 1 the exponent in (52) at order $j=m+1\le k$ is
$d-m=k-1+\alpha-m=k+\alpha-j\ge\alpha>0$, so the derivatives of order at most
$k$ tend to $0$ at a singular point; this lemma is what turns that into
differentiability there, and it also supplies the continuity at the singular
point which the hypothesis ``$F$ continuous on $U$'' of Lemma 2 of the paper
(paper.tex:538--539) silently presumes.

The hypothesis on $Z$ is the paper's own.  At paper.tex:156 the complex
gradient $f$ is asserted to be continuous and, in the same sentence, its zero
set $f^{-1}(0)$ to be discrete; a discrete subset that is closed in $\Omega$ ---
which $f^{-1}(0)$ is, by that continuity --- meets every compact subset of
$\Omega$ in a finite set.  So local finiteness is what the paper supplies, and
it is also what the argument below needs: mere isolation of each point of $Z$
within $Z$ would leave $\Omega-Z$ possibly non-open, since such a $Z$ may
accumulate at a point of $\Omega$ outside $Z$, and steps (b) and (d) both read
$\Omega-Z$ as an open set.
\end{lemma}
\begin{proof}
\emph{Step 0: two consequences of the hypothesis on $Z$.}  Both are used below,
and both come from the local finiteness of $Z$.  Fix the following construction,
available at any $\zeta\in\Omega$: since $\Omega$ is open there is $\rho>0$ with
$\kappa_\zeta=\{z:|z-\zeta|\le\rho\}\subseteq\Omega$, and $\kappa_\zeta$ is
compact, so $Z\cap\kappa_\zeta$ is finite by hypothesis.

\emph{(0a) $\Omega-Z$ is open.}  Let $z\in\Omega-Z$ and form $\kappa_z$ and
$\rho$ as above.  The set $Z\cap\kappa_z$ is finite and does not contain $z$,
so each of its finitely many elements $\zeta$ has $|\zeta-z|>0$; let
$\varepsilon$ be the least of those distances if $Z\cap\kappa_z\neq\emptyset$,
and $\varepsilon=\rho$ otherwise, and put $s=\min\{\varepsilon,\rho\}>0$.  If
$|z'-z|<s$ then $z'\in\kappa_z\subseteq\Omega$, and $z'\notin Z$: otherwise
$z'\in Z\cap\kappa_z$ with $|z'-z|<s\le\varepsilon\le|z'-z|$, a contradiction.
So $\{z':|z'-z|<s\}\subseteq\Omega-Z$, and $z$ is interior to $\Omega-Z$.

\emph{(0b) Each point of $Z$ is isolated in $Z$.}  Let $z_0\in Z$ and form
$\kappa_{z_0}$ and $\rho$ as above.  The set $Z\cap\kappa_{z_0}$ is finite, so
its elements other than $z_0$ are finitely many and each has $|\zeta-z_0|>0$;
let $\varepsilon$ be the least of those distances if there are any, and
$\varepsilon=\rho$ otherwise, and put $r_0=\min\{\varepsilon,\rho\}>0$.  If
$z\in Z$ with $|z-z_0|<r_0$ then $z\in\kappa_{z_0}$, so $z\in Z\cap\kappa_{z_0}$,
and $z\neq z_0$ would give $\varepsilon\le|z-z_0|<r_0\le\varepsilon$; hence
$z=z_0$.  That is, $\{z:|z-z_0|<r_0\}\cap Z=\{z_0\}$.

We prove by induction on $i$ that $w$ is of class $C^{i}$ on $\Omega$ for every
$i$ with $1\le i\le k$; the case $i=k$ is the assertion.

\emph{Base case $i=1$.}  This is (H1).

\emph{Inductive step.}  Let $1\le i<k$ and assume $w$ is of class $C^{i}$ on
$\Omega$.  We show that $w$ is of class $C^{i+1}$ on $\Omega$.  It suffices to
show that for every word $v'$ of length $i$ the function
\[
  g=D^{v'}w
\]
is of class $C^{1}$ on $\Omega$: a function of class $C^{i}$ on an open set all
of whose iterated partial derivatives of order $i$ are of class $C^{1}$ there
is of class $C^{i+1}$ there, since the derivative of order $i+1$ is assembled
from the first derivatives of the derivatives of order $i$, and these are then
continuous.

Fix such a $v'$.  We check that $g$ is differentiable at every point of
$\Omega$ and that $\mathrm{d}g$ is continuous on $\Omega$.

\emph{(a) $g$ is continuous on $\Omega$.}  $w$ is $C^{i}$ on $\Omega$ by the
inductive hypothesis and $|v'|=i$, so $g=D^{v'}w$ is continuous on $\Omega$:
the iterated derivative of order $i$ is continuous there, and $g$ is its value
on the fixed tuple of basis vectors prescribed by $v'$, by
\noderef{PartialWord}.

\emph{(b) $g$ is $C^\infty$, in particular $C^1$, on $\Omega-Z$.}  The set
$\Omega-Z$ is open by (0a), and by (H2) $w$ is $C^\infty$ on it; hence so is
each of its iterated partial derivatives, in particular $g$.  This is where the
openness of $\Omega-Z$ is needed: on a non-open set the smoothness of (H2)
would be smoothness within the set and would not give differentiability at a
point of $\Omega$.

\emph{(c) At each point of $Z$.}  Let $z_0\in Z$ and take $\delta>0$ and $B_j$
as in (H3) for this $z_0$.  Choose $r>0$ with
\[
  r<\delta,\qquad \{z:|z-z_0|<r\}\subseteq\Omega,\qquad
  \{z:|z-z_0|<r\}\cap Z=\{z_0\} ,
\]
which is possible: the first two by $\delta>0$ and openness of $\Omega$ at
$z_0\in Z\subseteq\Omega$, and the third by (0b), after shrinking $r$ below the
$r_0$ it provides.  Put
$\Omega'=\{z:|z-z_0|<r\}$, an open set containing $z_0$ with
$\Omega'\subseteq\Omega$ and $\Omega'-\{z_0\}\subseteq\Omega-Z$.

By (a), $g$ is continuous on $\Omega'$.  By (b), $g$ is differentiable at every
point of $\Omega'-\{z_0\}$.  For the derivative limit, let
$z\in\Omega'-\{z_0\}$.  The two first partial derivatives of $g$ at $z$ are
$D^{[0]+\!\!+v'}w(z)$ and $D^{[1]+\!\!+v'}w(z)$: by the recursion of
\noderef{PartialWord} one has $D^{\,\iota::v'}=\partial_\iota\circ D^{v'}$, so
prefixing a letter to the word is applying that partial derivative last.  And
an $\mathbb{R}$-linear map $\mathbb{C}\to\mathbb{C}$ has operator norm at most
the sum of the moduli of its values on $1$ and $i$.  Hence
\[
  \bigl\|\mathrm{d}g(z)\bigr\|
  \le\bigl|D^{[0]+\!\!+v'}w(z)\bigr|+\bigl|D^{[1]+\!\!+v'}w(z)\bigr|
  \le\sum_{|v|=i+1}\bigl|D^{v}w(z)\bigr|
  \le B_{i+1}\,|z-z_0|^{\,k+\alpha-(i+1)},
\]
the middle inequality because the two words $[0]\!+\!\!+v'$ and
$[1]\!+\!\!+v'$ are among the words of length $i+1$ and all terms of the sum
are nonnegative, and the last by (H3) applied with $j=i+1$, which is legitimate
since $1\le i+1\le k$ and $0<|z-z_0|<r<\delta$.

Now $i+1\le k$ gives $k+\alpha-(i+1)\ge\alpha>0$, so
$|z-z_0|^{\,k+\alpha-(i+1)}\to0$ as $z\to z_0$.  Also $B_{i+1}\ge0$, since the
left-hand side of (H3) is a sum of moduli and the power of $|z-z_0|$ is
positive for at least one admissible $z$.  Therefore
\[
  \mathrm{d}g(z)\ \longrightarrow\ 0\qquad
  \text{as }z\to z_0\text{ within }\Omega'-\{z_0\}.
\]
The hypotheses of \noderef{FDerivAcrossPoint} now hold with the open set
$\Omega'$, the point $z_0\in\Omega'$, the function $g$, and $L=0$: $g$ is
continuous on $\Omega'$, differentiable at every point of $\Omega'-\{z_0\}$,
and its derivative tends to $0$ at $z_0$ within $\Omega'-\{z_0\}$.  It
therefore gives that $g$ is differentiable at $z_0$ with $\mathrm{d}g(z_0)=0$,
and that $\mathrm{d}g$ is continuous at $z_0$ within $\Omega'$.  Since
$\Omega'$ is an open neighbourhood of $z_0$, continuity within $\Omega'$ at
$z_0$ is continuity at $z_0$, hence also continuity at $z_0$ within $\Omega$.

\emph{(d) Conclusion for $g$.}  By (b), $g$ is differentiable at every point of
$\Omega-Z$ and $\mathrm{d}g$ is continuous there; by (c), $g$ is differentiable
at every point of $Z$ and $\mathrm{d}g$ is continuous at each such point within
$\Omega$.  As $\Omega=(\Omega-Z)\cup Z$, the function $g$ is differentiable
everywhere on $\Omega$ and $\mathrm{d}g$ is continuous on $\Omega$; that is, $g$
is of class $C^{1}$ on $\Omega$.

Since $v'$ was an arbitrary word of length $i$, the remark opening the
inductive step gives that $w$ is of class $C^{i+1}$ on $\Omega$, completing the
induction and the proof.
\end{proof}
